% Lösungen Lineare Gleichungssysteme (zusammenbau v0.4, Heft)
\einheitenkopf{Lineare Gleichungssysteme}

\erg{25}{a) von beiden \quad b) $y$: $5$, $4$, $3$, $2$ \quad c) $y = -2x + 6$}
\erg{26}{a) $x = 6$ Zweierzelte, $y = 8$ Viererzelte; Probe: $6 + 8 = 14$, $12 + 32 = 44$. Keine andere Möglichkeit: Tauscht man ein Zweierzelt gegen ein Viererzelt, kommen $2$ Plätze dazu – die Platzzahl ändert sich bei jedem Tausch, nur eine Aufteilung ergibt $44$. \quad b) $x = 8$ Vierertische, $y = 7$ Sechsertische; Probe: $8 + 7 = 15$, $32 + 42 = 74$. Keine andere Möglichkeit: Die Geraden zu $x + y = 15$ und $4x + 6y = 74$ haben verschiedene Steigungen und schneiden sich in genau einem Punkt.}
\erg{27}{a) Gleichung II \quad b) $x = 3$, $y = 7$; Lösung $(3|7)$ \quad c) $x = 5$, $y = 9$; Lösung $(5|9)$}
\erg{28}{a) $y = 54 - x$ in II: $7x + 216 - 4x = 316{,}20$, $x = 33{,}40$; $y = 20{,}60$. Ein Helm kostet $33{,}40\,€$, ein Schloss $20{,}60\,€$. \quad b) $y = 48 - x$ in II: $5x + 96 - 2x = 199{,}20$, $x = 34{,}40$; $y = 13{,}60$. Ein Kanu kostet $34{,}40\,€$, ein Kajak $13{,}60\,€$.}
\erg{29}{a) rechts: $3{,}10\,€$ und $3{,}70\,€$; vor $x$ und $y$: $3$, $2$ und $1$, $4$; $x$ Preis eines Brötchens, $y$ Preis einer Brezel (in €) \quad b) I: $x + y = 4{,}60$, II: $3x + 4y = 15{,}80$ \quad c) I: $x + y = 50$, II: $2{,}80x + 4{,}50y = 177{,}40$}
\erg{30}{a) I: $x + y = 4{,}10$, II: $5x + 2y = 13{,}60$ \quad b) I: $x + y = 23{,}50$, II: $2x + 3y = 56{,}50$ \quad c) II: Mia kauft zusammen $14$ Äpfel und Mangos. $a = 14 - m$ in I: $8{,}40 + 0{,}80m = 11{,}60$, $m = 4$, $a = 10$: $10$ Äpfel und $4$ Mangos. \quad d) II: Der Verein kauft zusammen $14$ Bälle und Hütchen. $h = 14 - b$ in I: $11b + 49 = 104$, $b = 5$, $h = 9$: $5$ Bälle und $9$ Hütchen. \quad e) $x$: Preis eines Zelts in €, $y$: Preis eines Schlafsacks in € (nicht die Anzahlen) \quad f) $x$: Preis einer Tischtennisplatte in €, $y$: Preis eines Kickers in € (nicht die Anzahlen) \quad g) $x$ Preis einer Gitarrenstunde, $y$ Preis einer Klavierstunde: I: $2x + 3y = 94{,}50$, II: $x + 2y = 57{,}00$; $x = 57 - 2y$ in I: $114 - y = 94{,}50$, $y = 19{,}50$; $x = 18{,}00$. Eine Gitarrenstunde kostet $18{,}00\,€$, eine Klavierstunde $19{,}50\,€$. \quad h) $x$ Preis einer Monatskarte, $y$ Preis einer Kurskarte: I: $3x + 4y = 177$, II: $x + 2y = 71$; $x = 71 - 2y$ in I: $213 - 2y = 177$, $y = 18$; $x = 35$. Eine Monatskarte kostet $35\,€$, eine Kurskarte $18\,€$.}
